\part{その他}
	\chapter{無限級数公式集}
		\section{\texorpdfstring{$\sum_{k=1}^{\infty}{\frac{\sin{k}}{k}} = \frac{\pi-1}{2}$}{Σ\_\{k=1\}\string^∞ sin(k)/k = (π-1)/2}}
			\begin{proof}
				\quad\par
				下図のような鋸波$y=f(x)$をフーリエ級数展開して求める。
				\begin{figure}[H]
					\centering
					\includegraphics[keepaspectratio, scale=0.5]
					{parts/complexAnalysis/imgs/series_frac{sin(n)}{n}/pic.pdf}
					\caption{鋸波}
				\end{figure}
				$f(x)$は奇関数なので$\cos$成分はゼロ、つまり$\forall n, a_n=0$。$\sin(nx)$の係数は
				\[b_n=\frac{1}{\pi}\integrate{-\pi}{\pi}{x\sin(nx)}{}{x}=\frac{2}{\pi}\integrate{-\pi}{0}{x\sin(nx)}{}{x}=(-1)^{n-1}\frac{2}{n}\]
				よって
				\[x\sim2\left(\frac{\sin x}{1}-\frac{\sin 2x}{2}+\frac{\sin 3x}{3}-\dotsb\right)\]
				$x$を$x-\pi$で置き換えると
				\[x-\pi \sim -2\left(\frac{\sin x}{1}+\frac{\sin2x}{2}+\frac{\sin3x}{3}+\dotsb\right)\]
				特に$x$=1とすれば公式を得る。
			\end{proof}
		\section{\texorpdfstring{$\sum_{n=1}^\infty{\frac{1}{1+n^2}}=\frac{-1}{2}+\frac{\pi}{2}\coth{\pi}$}{Σ\_\{n=1\}\string^∞ 1/(1+n\string^2) = -1/2 + (π/2)coth(π)}}
			大阪大学 大学院 工学研究科 平成21年度 院試 数学 問題4 の解答を参照。
		\section{\texorpdfstring{$\sum_{n=1}^\infty{\frac{1}{(1+n^2)^2}}=\frac{-1}{2}+\frac{\pi^2}{4\sinh^2{\pi}}+\frac{\pi}{4}\coth\pi$}{Σ\_\{n=1\}\string^∞ 1/(1+n\string^2)\string^2 = -1/2 + π\string^2/(4sinh\string^2(π)) + (π/4)coth(π)}}
			大阪大学 大学院 工学研究科 平成20年度 院試 数学 問題4 の解答を参照。